\section{Future directions}
\label{sec:future}
The campaign localises two distinct limits with concrete, separable levers. The discovery frontier
is bounded by \emph{vetting resolution} (model improvements saturate at DECaLS $0.262''$/px), and
the footprint is bounded by \emph{training instrument} (v3 is DECam-specific). The priorities follow
directly.

\paragraph{1. Euclid pre-screening (the decisive discovery lever).} The cross-validation
(\S\ref{sec:v3}) shows the productive mode: v3 pre-screens the DESI imaging, Euclid confirms at
$0.1''$. Q1's $\sim$63\,deg$^2$ over three Deep Fields overlaps only $\sim$0.5\% of a DESI-south
sweep, which is why net-new yield there is bounded; \emph{Euclid DR1} (wide-area) inverts this. The
ready action now is a targeted sweep restricted to the Euclid-overlap sky, where \emph{every}
candidate is escalatable --- the one configuration in which v3's separation and the resolution-
breaking vet both apply to every object.

\paragraph{2. BASS/MzLS-north retraining.} v3 does not transfer to the northern third of the sky
(D3). Mining a BASS/MzLS-native lens-mimic bank and retraining the members (the v3 recipe is
instrument-agnostic; only the negatives change) would unlock DR9/DR11-north and roughly triple the
addressable footprint.

\paragraph{3. \emph{i}-band / foundation-model revival for small-$\theta_E$.} DR10/DR11-south are
native \emph{griz}. The deferred \aion{}/$i$-band lever (A4) targets the red-lens/red-companion
degeneracy and the small-$\theta_E$ bin specifically; it is worth re-testing under the narrower
lens-vs-mimic task (rather than the random-FPR task where it failed in v2) on the native-$i$ rows
the bank already carries.

\paragraph{4. A pixel-level lens-vs-mimic head.} The learned logistic head on member scores beats
the average at $\phi{=}0.05$ but overfits the strict tail and does not generalise to the independent
Euclid ranking. A small CNN head trained directly on pixels with positives$=$lenses,
negatives$=$the typed mimic bank (rather than on the eight member scores) is the natural next step,
used as a post-conformal re-ranker.

\paragraph{5. The active-learning loop.} The \texttt{lensjudge} harness can export confident-D
rejections (true mimics, \emph{not} objects that flipped to A/B at higher resolution) as fresh,
typed hard negatives for the next-generation bank --- closing a mine$\to$grade$\to$retrain loop that
continually sharpens the dominant contaminant classes.

\paragraph{6. A full DR11-south candidate catalogue.} \emph{Realised --- see \S\ref{sec:v4}.} v3blend8
runs on DR11-south with essentially no adaptation (the pipeline is release-parametrised); v4 went
further --- replacing the per-release threshold recalibration with a calibrated-mean selector and adding
a release-native warm-start fine-tune --- and re-swept the full footprint into the v4 candidate set. The
honest caveat stands: net-new yield is resolution-bounded, so the v4 catalogue is recall-richer but its
new candidates remain pending higher-resolution vetting.

\paragraph{Deployment recommendation.} For DR11-south, ship the \textbf{v4} release-native fine-tuned
ensemble with the \textbf{calibrated-mean stage-1 selector} (\S\ref{sec:v4}); retain \textbf{v3blend8}
for DR10-south and the Euclid cross-validation. In all cases use the score \emph{mean} (not the learned
head, which does not generalise past $\phi{=}0.05$) and the LensJudge cascade (direct triage $+$ agentic
mimic adjudication $+$ HSC/Euclid escalation by coverage) for vetting. This keeps the v2-lean lens
recall, adds v3's contaminant-rejection edge and v4's release-native gain on the hard residual, and
routes the resolution-breaking confirmation to exactly the candidates with higher-resolution coverage.
